\documentclass[10pt,a4paper]{amsart}
\usepackage{amsmath,amssymb,amsthm,graphicx}
% Retained preamble for the Stacks Project chapter {CH} (shared project preamble).
% The chapter file pinned-original-source.tex only carries \input{{preamble}} on its first line,
% so no preamble text lives inside the chapter itself.  The Stacks Project uses one single shared
% preamble for every chapter: the snapshot stages it three times (sources/stacks-chow/preamble.tex,
% sources/stacks-etale-cohomology/preamble.tex and focused-seeds/AGP/sources/stacks-snapshot-20261011/
% preamble.tex) and those copies are identical in size.  The block below is copied mechanically,
% line for line, from sources/stacks-chow/preamble.tex lines 14-18 and 154-217: the Xy-pic setup and
% the native theorem declarations together with the mathematical macros.  Nothing was added, removed
% or redefined; the wrapper class and the amsmath/amssymb/amsthm/graphicx base packages are supplied
% by the assembler.
% For commutative diagrams we use Xy-pic
% Chapter: stacks-cohomology.tex -- Cohomology of Sheaves.
\usepackage[all]{xy}

% We use 2cell for 2-commutative diagrams.
\xyoption{2cell}

% Theorem environments.
%
\theoremstyle{plain}
\newtheorem{theorem}[subsection]{Theorem}
\newtheorem{proposition}[subsection]{Proposition}
\newtheorem{lemma}[subsection]{Lemma}

\theoremstyle{definition}
\newtheorem{definition}[subsection]{Definition}
\newtheorem{example}[subsection]{Example}
\newtheorem{exercise}[subsection]{Exercise}
\newtheorem{situation}[subsection]{Situation}

\theoremstyle{remark}
\newtheorem{remark}[subsection]{Remark}
\newtheorem{remarks}[subsection]{Remarks}

\numberwithin{equation}{subsection}

% Macros
%
\def\lim{\mathop{\mathrm{lim}}\nolimits}
\def\colim{\mathop{\mathrm{colim}}\nolimits}
\def\Spec{\mathop{\mathrm{Spec}}}
\def\Hom{\mathop{\mathrm{Hom}}\nolimits}
\def\Ext{\mathop{\mathrm{Ext}}\nolimits}
\def\SheafHom{\mathop{\mathcal{H}\!\mathit{om}}\nolimits}
\def\SheafExt{\mathop{\mathcal{E}\!\mathit{xt}}\nolimits}
\def\Sch{\mathit{Sch}}
\def\Mor{\mathop{\mathrm{Mor}}\nolimits}
\def\Ob{\mathop{\mathrm{Ob}}\nolimits}
\def\Sh{\mathop{\mathit{Sh}}\nolimits}
\def\NL{\mathop{N\!L}\nolimits}
\def\CH{\mathop{\mathrm{CH}}\nolimits}
\def\proetale{{pro\text{-}\acute{e}tale}}
\def\etale{{\acute{e}tale}}
\def\QCoh{\mathit{QCoh}}
\def\Ker{\mathop{\mathrm{Ker}}}
\def\Im{\mathop{\mathrm{Im}}}
\def\Coker{\mathop{\mathrm{Coker}}}
\def\Coim{\mathop{\mathrm{Coim}}}

% Boxtimes
%
\DeclareMathSymbol{\boxtimes}{\mathbin}{AMSa}{"02}

%
% Macros for moduli stacks/spaces
%
\def\QCohstack{\mathcal{QC}\!\mathit{oh}}
\def\Cohstack{\mathcal{C}\!\mathit{oh}}
\def\Spacesstack{\mathcal{S}\!\mathit{paces}}
\def\Quotfunctor{\mathrm{Quot}}
\def\Hilbfunctor{\mathrm{Hilb}}
\def\Curvesstack{\mathcal{C}\!\mathit{urves}}
\def\Polarizedstack{\mathcal{P}\!\mathit{olarized}}
\def\Complexesstack{\mathcal{C}\!\mathit{omplexes}}
% \Pic is the operator that assigns to X its picard group, usage \Pic(X)
% \Picardstack_{X/B} denotes the Picard stack of X over B
% \Picardfunctor_{X/B} denotes the Picard functor of X over B
\def\Pic{\mathop{\mathrm{Pic}}\nolimits}
\def\Picardstack{\mathcal{P}\!\mathit{ic}}
\def\Picardfunctor{\mathrm{Pic}}
\def\Deformationcategory{\mathcal{D}\!\mathit{ef}}
\usepackage{cleveref}
\crefname{originallemma}{original lemma}{original lemmas}
\Crefname{originallemma}{Original lemma}{Original lemmas}
\makeatletter
\newlabel{lemma-describe-higher-direct-images}{{01E4}{source}{}{sourcepointer-1}{}}
\newlabel{lemma-describe-higher-direct-images@cref}{{[originallemma][0][]01E4}{[1][0][]source}{}{}{}}
\makeatother
\crefname{originalremark}{original remark}{original remarks}
\Crefname{originalremark}{Original remark}{Original remarks}
\makeatletter
\newlabel{remark-elucidate-lemma}{{01F0}{source}{}{sourcepointer-2}{}}
\newlabel{remark-elucidate-lemma@cref}{{[originalremark][0][]01F0}{[1][0][]source}{}{}{}}
\makeatother
\title{Cohomology and higher direct images of a module agree with those of the underlying abelian sheaf}
\author{The Stacks Project Authors}
\date{stacks-project:cohomology.tex:local-snapshot-20261011; GNU FDL 1.2 or later}
\begin{document}
\maketitle
\noindent\textit{Editorial note.} All delivered mathematical text is the author's own, copied from the pinned Stacks Project chapter file cohomology.tex; no statement or proof step was rewritten, reordered or removed.\par\medskip
\noindent\textit{Editorial note.} The proof is complete as written and covers the bounded below complex case by the same argument; no step is referred to the reader.\par\medskip
\noindent\textit{Editorial note.} The counted spans carry 2 navigation stubs and no bibliography citation.\par\medskip
\noindent\textit{Editorial note.} The preamble is the chapter's retained minimal preamble (sources/stacks-cohomology/retained-preamble.tex lines 1-81), a verbatim excerpt of the native preamble.\par\medskip
\noindent\textit{Editorial note.} The label lemma-modules-abelian also occurs in etale-cohomology.tex. There the statement compares H\textasciicircum{}i(G, -) for modules over a group ring with H\textasciicircum{}i(G, -) for modules over the group, which is a different proposition; only the label coincides, so this card counts its own statement.\par\medskip
\noindent\textit{Allowed standard prerequisites (editorial).}\begin{itemize}
\item Compatibility of derived functors along a morphism of ringed spaces. Same chapter, Remark remark-elucidate-lemma, invoked at proof lines 2202-2204. Supplies that one injective resolution computes both derived functors.
\item Description of higher direct images by derived pushforward. Same chapter, Lemma lemma-describe-higher-direct-images, invoked at proof lines 2208-2209. Supplies part (2) from part (1).
\end{itemize}\medskip
\ifdefined\hypertarget\hypertarget{sourcepointer-1}{}\fi
\noindent\textit{Source reference.} Original lemma 01E4: sources/stacks-cohomology/pinned-original-source.tex, source lines 592--592. Pointer outside the counted span, invoked once for context; the load-bearing argument of the counted statement is the author's proof reproduced in full in the proof span.\par\medskip
\ifdefined\hypertarget\hypertarget{sourcepointer-2}{}\fi
\noindent\textit{Source reference.} Original remark 01F0: sources/stacks-cohomology/pinned-original-source.tex, source lines 2147--2147. Pointer outside the counted span, invoked once for context; the load-bearing argument of the counted statement is the author's proof reproduced in full in the proof span.\par\medskip
\begin{lemma}
\label{lemma-modules-abelian}
Let $X$ be a ringed space.
Let $\mathcal{F}$ be an $\mathcal{O}_X$-module.
\begin{enumerate}
\item The cohomology groups $H^i(U, \mathcal{F})$ for $U \subset X$ open
of $\mathcal{F}$ computed as an $\mathcal{O}_X$-module, or computed as an
abelian sheaf are identical.
\item Let $f : X \to Y$ be a morphism of ringed spaces.
The higher direct images $R^if_*\mathcal{F}$ of $\mathcal{F}$
computed as an $\mathcal{O}_X$-module, or computed as an abelian sheaf
are identical.
\end{enumerate}
There are similar statements in the case of bounded below
complexes of $\mathcal{O}_X$-modules.
\end{lemma}

\begin{proof}
Consider the morphism of ringed spaces
$(X, \mathcal{O}_X) \to (X, \underline{\mathbf{Z}}_X)$ given
by the identity on the underlying topological space and by
the unique map of sheaves of rings
$\underline{\mathbf{Z}}_X \to \mathcal{O}_X$.
Let $\mathcal{F}$ be an $\mathcal{O}_X$-module.
Denote $\mathcal{F}_{ab}$ the same sheaf seen as an
$\underline{\mathbf{Z}}_X$-module, i.e., seen as a sheaf of
abelian groups. Let
$\mathcal{F} \to \mathcal{I}^\bullet$ be an injective resolution.
By Remark \ref{remark-elucidate-lemma} we see that
$\Gamma(X, \mathcal{I}^\bullet)$ computes both
$R\Gamma(X, \mathcal{F})$ and $R\Gamma(X, \mathcal{F}_{ab})$.
This proves (1).

\medskip\noindent
To prove (2) we use (1) and Lemma \ref{lemma-describe-higher-direct-images}.
The result follows immediately.
\end{proof}

\end{document}
